\documentclass[10pt,a4paper]{article}

\usepackage[utf8]{inputenc}
\usepackage[T1]{fontenc}
\usepackage[default]{sourcesanspro}   % clean sans-serif (IBM Plex Sans feel)
\input{glyphtounicode}
\pdfgentounicode=1
\usepackage[top=0.42in,bottom=0.42in,left=0.46in,right=0.46in]{geometry}
\usepackage[colorlinks=true,urlcolor=cvaccent,linkcolor=cvaccent]{hyperref}
\usepackage{enumitem}
\usepackage[dvipsnames]{xcolor}

% ── Palette (derived from anderson-ferreira-83.github.io) ──────────────
%   site: bg #020406 · accent #53c8ff · text #f3f7fb · muted #9aafc2
%   cv (white paper): darken accent for legibility, keep brand family
\definecolor{cvdark}{HTML}{081420}    % near-black navy  → headings / name
\definecolor{cvaccent}{HTML}{0A7FB5}  % darkened cyan    → links, dates, rules
\definecolor{cvmuted}{HTML}{456272}   % blue-grey        → institutions, labels
\definecolor{cvrule}{HTML}{B8D4E8}    % light cyan       → decorative hairlines
\definecolor{cvbullet}{HTML}{1A9ED4}  % mid-cyan         → bullet dots

\pagestyle{empty}
\setlength{\parindent}{0pt}
\setlength{\parskip}{0pt}
\setlist[itemize]{%
  leftmargin=1.2em,
  itemsep=1pt,
  topsep=2pt,
  parsep=0pt,
  label={\small\color{cvbullet}\textbullet}
}
\urlstyle{same}

% ── Commands ───────────────────────────────────────────────────────────

% Section with numbered kicker  →  01 · Título
\newcommand{\cvsection}[2]{%
  \vspace{0.42em}
  {\footnotesize\bfseries\color{cvmuted}#1}\enspace
  {\large\bfseries\color{cvdark}#2}\par
  \vspace{0.04em}
  {\color{cvrule}\hrule height 0.5pt}
  \vspace{0.18em}
}

% Experience / Project entry: {Título}{Período}{Subtítulo}{Corpo}
\newcommand{\entry}[4]{%
  \noindent\textbf{\color{cvdark}#1}\hfill{\small\bfseries\color{cvaccent}#2}\par
  {\small\itshape\color{cvmuted}#3}\par
  \vspace{0.04em}
  {\small #4}
  \vspace{0.12em}
}

% Education: {Grau e Curso}{Instituição}{Período}
\newcommand{\cvdegree}[3]{%
  \noindent\textbf{\color{cvdark}#1}\hfill{\small\color{cvaccent}#3}\par
  {\small\itshape\color{cvmuted}#2}\par
  \vspace{0.06em}
}

% Skill row: {Categoria}{Habilidades}
\newcommand{\skill}[2]{%
  \noindent\makebox[2.9cm][l]{{\small\bfseries\color{cvdark}#1}}{\small\color{cvmuted}#2}\par
  \vspace{0.03em}
}

\begin{document}

% ── Top accent stripe (mirrors featured-card stripe on the site) ────────
{\color{cvaccent}\noindent\rule{\linewidth}{1.8pt}}\par
\vspace{0.22em}

% ── Name block ──────────────────────────────────────────────────────────
{\fontsize{19.5}{23}\selectfont\bfseries\color{cvdark} Anderson H. R. Ferreira}\par
\vspace{0.08em}
{\footnotesize\bfseries\color{cvmuted}\MakeUppercase{From Signals to Intelligence}}\par
\vspace{0.05em}
{\small\color{cvaccent}%
  Dados \& Analytics · Machine Learning · Processamento de Sinais · IoT%
  \enspace|\enspace Doutorando UNICAMP%
}\par
\vspace{0.18em}

% ── Contact row ─────────────────────────────────────────────────────────
{\small\color{cvmuted}%
  +55\,11\,95850-9690\enspace·\enspace
  \href{mailto:a058899@dac.unicamp.br}{a058899@dac.unicamp.br}\enspace·\enspace
  Salto, SP\enspace·\enspace
  \href{https://anderson-ferreira-83.github.io/}{anderson-ferreira-83.github.io}\enspace·\enspace
  \href{https://www.linkedin.com/in/anderson-ferreira-1a473138/?locale=pt}{LinkedIn}\enspace·\enspace
  \href{https://github.com/anderson-ferreira-83}{GitHub}\\[2pt]
  Defesa prevista: dez/2026\enspace·\enspace
  Disponibilidade híbrida em São Paulo: jun/2026--set/2026%
}\par
\vspace{0.12em}
{\color{cvrule}\noindent\rule{\linewidth}{0.4pt}}\par
\vspace{0.18em}

% ── 00 · Resumo ─────────────────────────────────────────────────────────
\cvsection{00}{Resumo}
{\small Doutorando em Engenharia Mecânica na UNICAMP, com atuação aplicada em dados, machine
learning, cloud e estatística. Experiência em modelagem preditiva, validação de modelos, feature
engineering e implementação com Python, SQL, FastAPI, Oracle e AWS. Repertório construído em
projetos de IoT e saúde, com técnicas transferíveis para analytics, risco, fraude, churn e apoio
a decisão.}

% ── 01 · Competências ───────────────────────────────────────────────────
\cvsection{01}{Competências Técnicas}
\skill{Linguagens}{Python · SQL · Java · JavaScript}
\skill{Frameworks}{FastAPI · scikit-learn · Docker}
\skill{ML / Dados}{Random Forest · SMOTE · SHAP · ETL · EDA · Feature Engineering ·
  Cross-validation · Threshold Analysis · Estatística Aplicada}
\skill{Cloud / Infra}{AWS · Oracle SQL}

% ── 02 · Projetos ───────────────────────────────────────────────────────
\cvsection{02}{Projetos de Dados e ML}

\entry{Predição de Risco com SMOTE}{Projeto aplicado}%
  {Saúde | orientador e desenvolvedor}{%
  \begin{itemize}
    \item Predição de risco de hipertensão com 4.240 amostras e 12 variáveis clínicas.
    \item Split estratificado, SMOTE apenas no treino e validação cruzada 5-fold, evitando data leakage em problema desbalanceado.
    \item Avaliação com Recall, F2-score, AUC-ROC, thresholds clínicos e interpretabilidade via SHAP.
    \item Abordagem transferível para fraude, churn, credit scoring e cenários com classe rara.
  \end{itemize}
}

\entry{Pipeline IoT End-to-End}{Projeto aplicado}%
  {ESP32 · FastAPI · Oracle XE}{%
  \begin{itemize}
    \item Pipeline de coleta, tratamento e classificação para 7 estados operacionais: ESP32 + MPU6050 → FastAPI → Oracle XE.
    \item 104 features candidatas geradas; 16 finais selecionadas. Random Forest (200 árvores): 97,34\% ±\,0,63\% em CV e 94,24\% em holdout.
    \item Mais de 210 mil registros documentados; inferência em produção no browser sem servidor de ML dedicado.
  \end{itemize}
}

% ── 03 · Experiência ────────────────────────────────────────────────────
\cvsection{03}{Experiência Profissional}

\entry{Professor Adjunto}{fev/2025--Atual}%
  {Cruzeiro do Sul Educacional}{%
  \begin{itemize}
    \item Docência em POO (Java), Estrutura de Dados I e II, Programação Web e Bancos de Dados.
    \item Atividades práticas com Java, Python e SQL, conectando algoritmos, modelagem e aplicações reais.
    \item Acompanhamento de projetos com integração entre software, dados e bancos de dados.
  \end{itemize}
}

\entry{Pesquisador de Doutorado}{2018--Atual}%
  {UNICAMP --- Engenharia Mecânica}{%
  \begin{itemize}
    \item Pesquisa em processamento de sinais, modelagem estatística, sensores e classificação de estados operacionais.
    \item Publicação aceita no \textit{Mechanical Systems and Signal Processing} (MSSP), 2026 --- comparando 5 geometrias com abordagem orientada a dados.
  \end{itemize}
}

\entry{Professor Universitário}{2014--2021}%
  {CEUNSP}{%
  \begin{itemize}
    \item Disciplinas de Probabilidade e Estatística, Físicas Experimentais, Termodinâmica e Transferência de Calor.
    \item Atividades de regressão, testes de hipótese, inferência estatística e simulação computacional.
  \end{itemize}
}

% ── 04 · Formação ────────────────────────────────────────────────────────
\cvsection{04}{Formação}
\cvdegree{Doutorado em Engenharia Mecânica}{UNICAMP}{2018--2026 (em andamento)}
\cvdegree{CST em Ciência de Dados}{Estácio}{2025--2027 (em andamento)}
\cvdegree{Mestrado em Engenharia Mecânica}{UNICAMP}{2011--2013}
\cvdegree{Licenciatura em Física}{UNICAMP}{2006--2010}

% ── 05 · Certificações ───────────────────────────────────────────────────
\cvsection{05}{Certificações e Publicação}
{\small
  AWS Certified AI Practitioner\enspace·\enspace
  Applied Machine Learning in Python\enspace·\enspace
  IBM Data Engineering Professional Certificate\par
  \vspace{0.06em}
  \textit{Mechanical Systems and Signal Processing} (MSSP), 2026
  \enspace|\enspace Impact Factor \textbf{\color{cvaccent}8.9}
}

% ── Bottom hairline ───────────────────────────────────────────────────────
\vspace{0.34em}
{\color{cvrule}\noindent\rule{\linewidth}{0.4pt}}\par

\end{document}
